% !TEX program = xelatex
\documentclass[aspectratio=169,11pt]{beamer}
\usetheme{Madrid}
\usepackage{fontspec}
\setsansfont{DejaVu Sans}
\setmainfont{DejaVu Serif}
\usepackage{polyglossia}
\setdefaultlanguage{turkish}
\usepackage{amsmath,amssymb}
\usepackage{booktabs}
\usepackage{multicol}
\usepackage{xcolor}

% ---------- Renkler ----------
\definecolor{anamavi}{RGB}{0,45,98}
\definecolor{cozumyesil}{RGB}{0,110,50}
\definecolor{uyarikirmizi}{RGB}{180,20,20}
\setbeamercolor{structure}{fg=anamavi}
\setbeamercolor{frametitle}{bg=anamavi,fg=white}
\setbeamercolor{title}{bg=anamavi,fg=white}
\setbeamercolor{palette primary}{bg=anamavi,fg=white}
\setbeamercolor{palette secondary}{bg=anamavi!85!black,fg=white}
\setbeamercolor{palette tertiary}{bg=anamavi!70!black,fg=white}
\setbeamercolor{block title}{bg=anamavi,fg=white}
\setbeamercolor{block body}{bg=anamavi!7}
\setbeamercolor{block title example}{bg=cozumyesil,fg=white}
\setbeamercolor{block body example}{bg=cozumyesil!8}
\setbeamercolor{block title alerted}{bg=uyarikirmizi,fg=white}
\setbeamercolor{block body alerted}{bg=uyarikirmizi!7}
\setbeamertemplate{navigation symbols}{}
\setbeamertemplate{blocks}[rounded][shadow=false]

% ---------- Soru / Çözüm kalıbı ----------
% \soru{başlık}{soru metni}{çözüm metni}: önce soru, bir tıklamayla çözüm
\newcommand{\soru}[3]{%
\begin{frame}{#1}
  \begin{block}{Soru}\large #2\end{block}
  \pause
  \begin{exampleblock}{Çözüm}\large #3\end{exampleblock}
\end{frame}}
\newcommand{\ab}[1]{\overline{#1}}

\title[MMAT105 -- Hafta 12]{MMAT105 Genel Matematik}
\subtitle{Hafta 12: Oran ve Orantı (Doğru, Ters, Bileşik Orantı), Ortalamalar}
\author{}
\date{2026--2027 Güz Dönemi}

\begin{document}

% =====================================================
\begin{frame}
\titlepage
\end{frame}

% (iletişim çerçevesi kaldırıldı)

\begin{frame}{Geçen Hafta Hatırlatma}
\begin{columns}[t]
\column{0.58\textwidth}
\textbf{Hesaplayınız:}
\begin{enumerate}
\item \(x^2-9x+20\) ifadesini çarpanlarına ayırınız.
\item \(a+b=7\) ve \(a\cdot b=10\) ise \(a^2+b^2=?\)
\item \(\dfrac{51^2-49^2}{20}=?\)
\end{enumerate}
\column{0.36\textwidth}
\pause
\textbf{Cevaplar:}
\begin{enumerate}
\item \((x-4)(x-5)\)
\item \(7^2-2\cdot 10=29\)
\item \(\dfrac{2\cdot 100}{20}=10\)
\end{enumerate}
\end{columns}
\end{frame}

\begin{frame}{Bugün}
\begin{columns}[t]
\column{0.32\textwidth}
\begin{block}{Ders 1}
\begin{enumerate}
\item Oran (birimsiz)
\item Orantı, içler--dışlar
\item \(\frac{a}{b}=\frac{c}{d}=k\) özellikleri
\item \(a:b:c\) ve paylaştırma
\item Harita ölçeği
\end{enumerate}
\end{block}
\column{0.32\textwidth}
\begin{block}{Ders 2}
\begin{enumerate}\setcounter{enumi}{5}
\item Doğru orantı: \(y=kx\)
\item Ters orantı: \(x\cdot y=k\)
\item İşçi, hız, musluk
\item Ters orantılı paylaştırma
\item Bileşik orantı
\end{enumerate}
\end{block}
\column{0.32\textwidth}
\begin{block}{Ders 3}
\begin{enumerate}\setcounter{enumi}{10}
\item Aritmetik ortalama
\item Sayı ekleme / çıkarma
\item Ağırlıklı ortalama
\item Geometrik, harmonik ort.
\item Ardışık sayılar
\end{enumerate}
\end{block}
\end{columns}
\end{frame}

% =====================================================
\section{Ders 1: Oran ve Orantı}
\begin{frame}
\centering\Large\color{anamavi}\textbf{Ders 1}\\[6pt]\normalsize Oran, Orantı, Orantı Sabiti, Paylaştırma ve Ölçek
\end{frame}

\subsection{Oran ve orantı}
\begin{frame}{Oran ve Orantı}
\begin{block}{Tanım}
\begin{itemize}
\item \textbf{Oran:} İki çokluğun bölümü: \(\dfrac{a}{b}\) \ (\(b\neq 0\)). Aynı birimli çoklukların oranı \textbf{birimsizdir}; farklı birimlilerin oranı birimlidir (hız: km/sa).
\item \textbf{Orantı:} İki oranın eşitliği: \(\dfrac{a}{b}=\dfrac{c}{d}\) \ (\(b,d\neq 0\)).
\item \textbf{İçler--dışlar çarpımı:} \(\dfrac{a}{b}=\dfrac{c}{d}\ \Longleftrightarrow\ a\cdot d=b\cdot c\)
\end{itemize}
\end{block}
\begin{alertblock}{Sık yapılan hatalar}
\begin{itemize}
\item Birimler farklıysa önce eşitlenir: \(2\) m ile \(50\) cm'nin oranı \(\frac{2}{50}\) değil, \(\frac{200}{50}=4\).
\item \(\frac{a}{b}=\frac{2}{3}\), \(a=2\) ve \(b=3\) demek değildir: \(a=2k,\ b=3k\) (\(k\neq 0\)).
\end{itemize}
\end{alertblock}
\end{frame}

\soru{Oran -- Örnek 1}
{\(2\) m uzunluğundaki bir ipin \(50\) cm uzunluğundaki bir ipe oranı kaçtır? \(45\) dakikanın \(2\) saate oranı kaçtır?}
{Önce birimler eşitlenir:
\[\frac{2\text{ m}}{50\text{ cm}}=\frac{200\text{ cm}}{50\text{ cm}}=4,\qquad \frac{45\text{ dk}}{2\text{ sa}}=\frac{45\text{ dk}}{120\text{ dk}}=\frac{3}{8}\]
Sonuçlar birimsizdir.}

\soru{İçler--Dışlar -- Örnek 2}
{\(\dfrac{x-1}{x+3}=\dfrac{3}{5}\) olduğuna göre \(x\) kaçtır?}
{İçler--dışlar çarpımı:
\[5(x-1)=3(x+3)\ \Rightarrow\ 5x-5=3x+9\ \Rightarrow\ 2x=14\ \Rightarrow\ x=7\]
Kontrol: \(\dfrac{6}{10}=\dfrac{3}{5}\) \checkmark}

\soru{Oran -- Örnek 3}
{\(\dfrac{a}{b}=\dfrac{3}{4}\) olduğuna göre \(\dfrac{2a+b}{b-a}\) oranı kaçtır?}
{\(a=3k\), \(b=4k\) \ (\(k\neq 0\)) yazalım:
\[\frac{2a+b}{b-a}=\frac{6k+4k}{4k-3k}=\frac{10k}{k}=10\]}

\soru{İçler--Dışlar -- Örnek 4}
{\(\dfrac{a+b}{a-b}=\dfrac{5}{2}\) olduğuna göre \(\dfrac{a}{b}\) oranı kaçtır?}
{\[2(a+b)=5(a-b)\ \Rightarrow\ 2a+2b=5a-5b\ \Rightarrow\ 7b=3a\]
\[\frac{a}{b}=\frac{7}{3}\qquad(\text{kontrol: } a=7,\ b=3:\ \tfrac{10}{4}=\tfrac{5}{2})\]}

\subsection{Orantı sabiti}
\begin{frame}{Orantı Sabiti ve Özellikleri}
\begin{block}{\(\dfrac{a}{b}=\dfrac{c}{d}=\dfrac{e}{f}=k\) ise \ (paydalar sıfırdan farklı)}
\begin{itemize}
\item \(a=bk,\quad c=dk,\quad e=fk\)
\item \(\dfrac{a+c+e}{b+d+f}=k\) \qquad (payların toplamı / paydaların toplamı)
\item \(\dfrac{m\,a+n\,c}{m\,b+n\,d}=k\) \qquad (\(mb+nd\neq 0\))
\end{itemize}
\end{block}
\begin{block}{\(\dfrac{a}{2}=\dfrac{b}{3}=\dfrac{c}{5}=k\) yaklaşımı}
\(a=2k,\ b=3k,\ c=5k\) yazılır; verilen koşuldan \(k\) bulunur. Bu, \(a:b:c=2:3:5\) demektir.
\end{block}
\end{frame}

\soru{Orantı Sabiti -- Örnek 5}
{\(\dfrac{a}{2}=\dfrac{b}{3}=\dfrac{c}{5}\) ve \(a+b+c=40\) olduğuna göre \(a\), \(b\), \(c\) kaçtır?}
{Ortak değere \(k\) diyelim: \(a=2k,\ b=3k,\ c=5k\).
\[2k+3k+5k=10k=40\ \Rightarrow\ k=4\]
\[a=8,\qquad b=12,\qquad c=20\]}

\soru{Orantı Sabiti -- Örnek 6}
{\(\dfrac{x}{2}=\dfrac{y}{3}=\dfrac{z}{5}\) ve \(2x+y-z=10\) olduğuna göre \(x\) kaçtır?}
{\(x=2k,\ y=3k,\ z=5k\):
\[2(2k)+3k-5k=10\ \Rightarrow\ 2k=10\ \Rightarrow\ k=5\]
\[x=2\cdot 5=10\]}

\soru{Orantı Sabiti -- Örnek 7}
{\(\dfrac{a}{b}=\dfrac{c}{d}=\dfrac{e}{f}=\dfrac{3}{4}\) ve \(b+d+f=20\) ise \(a+c+e\) kaçtır? \ \(\dfrac{2a+5c}{2b+5d}\) kaçtır?}
{\[\frac{a+c+e}{b+d+f}=\frac{3}{4}\ \Rightarrow\ \frac{a+c+e}{20}=\frac{3}{4}\ \Rightarrow\ a+c+e=15\]
\(a=\tfrac34 b\), \(c=\tfrac34 d\) \(\Rightarrow\) \(2a+5c=\tfrac34(2b+5d)\) \(\Rightarrow\) \(\dfrac{2a+5c}{2b+5d}=\dfrac{3}{4}\).}

\soru{Oran Bulma -- Örnek 8}
{\(3a=4b=6c\) olduğuna göre \(a:b:c\) kaçtır? \ \(a+b+c=36\) ise \(a\) kaçtır?}
{Ortak değeri \(3,4,6\)'nın EKOK'u ile seçelim: \(3a=4b=6c=12k\).
\[a=4k,\quad b=3k,\quad c=2k\ \Rightarrow\ a:b:c=4:3:2\]
\(9k=36\Rightarrow k=4\Rightarrow a=16\).}

\soru{Oran Birleştirme -- Örnek 9}
{\(a:b=2:3\) ve \(b:c=4:5\) olduğuna göre \(a:b:c\) kaçtır? \ \(a+b+c=70\) ise \(c\) kaçtır?}
{Ortak terim \(b\) eşitlenir: \(\mathrm{EKOK}(3,4)=12\).
\[a:b=8:12,\qquad b:c=12:15\ \Rightarrow\ a:b:c=8:12:15\]
\(35k=70\Rightarrow k=2\Rightarrow c=15\cdot 2=30\).}

\subsection{Paylaştırma}
\soru{Paylaştırma -- Örnek 10}
{\(600\) TL, üç kişi arasında \(2\), \(3\) ve \(5\) sayıları ile orantılı olarak paylaştırılıyor. En çok ve en az alan kişi kaçar TL alır?}
{Paylar: \(2k,\ 3k,\ 5k\).
\[2k+3k+5k=10k=600\ \Rightarrow\ k=60\]
En çok: \(5\cdot 60=300\) TL; \quad en az: \(2\cdot 60=120\) TL.}

\subsection{Ölçek}
\begin{frame}{Harita Ölçeği}
\begin{block}{Tanım}
\[\text{Ölçek}=\frac{\text{Haritadaki uzunluk}}{\text{Gerçek uzunluk}}\qquad(\text{aynı birimle})\]
\(1:500\,000\) ölçeği: haritada \(1\) cm, gerçekte \(500\,000\) cm \(=5\) km'dir.
\end{block}
\begin{alertblock}{Sık yapılan hatalar}
\begin{itemize}
\item \(1\) km \(=1000\) m \(=100\,000\) cm. Birim dönüşümünü unutmayın.
\item Ölçek bir oran olduğundan birimsizdir: \(1:500\,000\).
\end{itemize}
\end{alertblock}
\end{frame}

\soru{Ölçek -- Örnek 11}
{\(1:500\,000\) ölçekli bir haritada iki şehir arası \(6\) cm'dir. Gerçek uzaklık kaç km'dir?}
{\[\text{Gerçek}=6\cdot 500\,000=3\,000\,000\text{ cm}\]
\(1\) km \(=100\,000\) cm \(\Rightarrow\) \(3\,000\,000\div 100\,000=30\) km.}

\soru{Ölçek -- Örnek 12}
{Gerçekte \(45\) km olan bir uzaklık haritada \(3\) cm ile gösteriliyor. Haritanın ölçeği nedir?}
{\(45\) km \(=45\cdot 100\,000=4\,500\,000\) cm.
\[\text{Ölçek}=\frac{3}{4\,500\,000}=\frac{1}{1\,500\,000}\ \Rightarrow\ 1:1\,500\,000\]}

\begin{frame}{Ders 1 -- Sınıfta Siz Çözün}
\begin{columns}[t]
\column{0.62\textwidth}
\begin{enumerate}
\item \(\dfrac{2x+1}{x+4}=\dfrac{5}{3}\) ise \(x=?\) \(\star\)
\item \(\dfrac{a}{3}=\dfrac{b}{4}=\dfrac{c}{6}\), \ \(a+b+c=52\) ise \(b=?\)
\item Bir üçgenin iç açıları \(2,3,4\) ile orantılı. En büyük açı kaç derece? \(\star\)
\item \(\dfrac{a}{b}=\dfrac{2}{5}\) ise \(\dfrac{3a+b}{b-a}=?\)
\item \(1:200\,000\) ölçekli haritada \(7{,}5\) cm kaç km'dir? \(\star\)
\end{enumerate}
\column{0.34\textwidth}
\pause
\textbf{Cevaplar:}
\begin{enumerate}
\item \(17\)
\item \(16\)
\item \(80^\circ\)
\item \(\dfrac{11}{3}\)
\item \(15\) km
\end{enumerate}
\end{columns}
\vspace{4pt}{\footnotesize \(\star\) işaretli soruların adım adım çözümü sonraki slaytlardadır.}
\end{frame}

% ---- Ders 1 alıştırma çözümleri ----
\soru{Ders 1 -- Alıştırma Çözümü (Soru 1)}
{\(\dfrac{2x+1}{x+4}=\dfrac{5}{3}\) olduğuna göre \(x\) kaçtır?}
{İçler--dışlar çarpımı:
\[3(2x+1)=5(x+4)\ \Rightarrow\ 6x+3=5x+20\ \Rightarrow\ x=17\]
Kontrol: \(\dfrac{35}{21}=\dfrac{5}{3}\) \checkmark}

\soru{Ders 1 -- Alıştırma Çözümü (Soru 3)}
{Bir üçgenin iç açıları \(2\), \(3\) ve \(4\) sayıları ile orantılıdır. En büyük açı kaç derecedir?}
{Açılar: \(2k,\ 3k,\ 4k\). Üçgenin iç açıları toplamı \(180^\circ\):
\[2k+3k+4k=9k=180^\circ\ \Rightarrow\ k=20^\circ\]
En büyük açı: \(4\cdot 20^\circ=80^\circ\) \quad (açılar: \(40^\circ,\ 60^\circ,\ 80^\circ\)).}

\soru{Ders 1 -- Alıştırma Çözümü (Soru 5)}
{\(1:200\,000\) ölçekli bir haritada \(7{,}5\) cm ile gösterilen uzaklık gerçekte kaç km'dir?}
{Haritada \(1\) cm \(\rightarrow\) gerçekte \(200\,000\) cm \(=2\) km.
\[7{,}5\cdot 200\,000=1\,500\,000\text{ cm}=15\text{ km}\]
(Kısa yol: \(7{,}5\cdot 2=15\) km.)}


% =====================================================
\section{Ders 2: Doğru, Ters ve Bileşik Orantı}
\begin{frame}
\centering\Large\color{anamavi}\textbf{Ders 2}\\[6pt]\normalsize Doğru Orantı, Ters Orantı ve Bileşik Orantı
\end{frame}

\subsection{Doğru orantı}
\begin{frame}{Doğru Orantı}
\begin{block}{Tanım}
\(x\) ve \(y\) pozitif çokluklar olmak üzere
\[\frac{y}{x}=k\quad\text{yani}\quad y=k\cdot x\qquad(k>0 \text{ sabit})\]
ise \(y\), \(x\) ile \textbf{doğru orantılıdır}: \(x\) kaç katına çıkarsa \(y\) de aynı kata çıkar.\\[2pt]
Örnek: alınan ürün miktarı -- ödenen para; harcanan benzin -- gidilen yol.
\end{block}
\begin{alertblock}{Sık yapılan hatalar}
``Biri artarken diğeri de artıyor'' yetmez; \textbf{oran sabit} olmalı.
\(y=x+3\) artandır ama doğru orantı değildir (\(\frac{4}{1}\neq\frac{5}{2}\)).
\end{alertblock}
\end{frame}

\soru{Doğru Orantı -- Örnek 13}
{\(y\), \(x\) ile doğru orantılıdır. \(x=4\) iken \(y=10\)'dur. \(x=10\) iken \(y\) kaçtır?}
{\[k=\frac{y}{x}=\frac{10}{4}=\frac{5}{2}\]
\[y=\frac{5}{2}\cdot 10=25\]
(Ya da: \(\dfrac{10}{4}=\dfrac{y}{10}\Rightarrow 4y=100\Rightarrow y=25\).)}

\soru{Doğru Orantı -- Örnek 14}
{\(x+2\) ile \(y-1\) doğru orantılıdır. \(x=4\) iken \(y=4\) oluyor. \(x=10\) iken \(y\) kaçtır?}
{Oran sabit: \(\dfrac{y-1}{x+2}=\dfrac{4-1}{4+2}=\dfrac{3}{6}=\dfrac{1}{2}\).\\[3pt]
\(x=10\): \ \(\dfrac{y-1}{12}=\dfrac{1}{2}\ \Rightarrow\ y-1=6\ \Rightarrow\ y=7\).}

\subsection{Ters orantı}
\begin{frame}{Ters Orantı}
\begin{block}{Tanım}
\(x\) ve \(y\) pozitif çokluklar olmak üzere
\[x\cdot y=k\qquad(k>0 \text{ sabit})\]
ise \(x\) ile \(y\) \textbf{ters orantılıdır}: \(x\) kaç katına çıkarsa \(y\) o kadar kata bölünür.\\[2pt]
Aynı iş için: işçi sayısı -- gün; hız -- zaman (aynı yol); musluk sayısı -- dolma süresi; dişli çarkta diş sayısı -- tur sayısı.
\end{block}
\begin{alertblock}{Sık yapılan hatalar}
Ters orantıda \textbf{çarpım} sabittir; oranı eşitlemeyin: \(x_1\cdot y_1=x_2\cdot y_2\).
\end{alertblock}
\end{frame}

\soru{Ters Orantı -- Örnek 15}
{\(x\) ile \(y\) ters orantılıdır. \(x=6\) iken \(y=8\)'dir. \(x=4\) iken \(y\) kaçtır?}
{\[x\cdot y=6\cdot 8=48\ \text{(sabit)}\]
\[4\cdot y=48\ \Rightarrow\ y=12\]
\(x\) azaldı, \(y\) arttı \checkmark}

\soru{İşçi -- Gün -- Örnek 16}
{\(4\) işçi bir duvarı \(15\) günde örüyor. Aynı nitelikteki \(10\) işçi aynı duvarı kaç günde örer?}
{İşçi sayısı artınca gün azalır \(\Rightarrow\) ters orantı.
\[4\cdot 15=10\cdot x\ \Rightarrow\ 60=10x\ \Rightarrow\ x=6\text{ gün}\]}

\soru{Dişli Çark -- Örnek 17}
{Birbirini çeviren iki dişli çarktan birinde \(20\), diğerinde \(30\) diş vardır. Küçük çark \(6\) tur attığında büyük çark kaç tur atar?}
{Diş sayısı \(\times\) tur sayısı sabittir (temas eden diş sayısı eşit):
\[20\cdot 6=30\cdot x\ \Rightarrow\ 120=30x\ \Rightarrow\ x=4\text{ tur}\]}

\soru{Hız -- Zaman -- Örnek 18}
{Bir araç bir yolu saatte \(60\) km hızla \(5\) saatte alıyor. Aynı yolu saatte \(75\) km hızla kaç saatte alır?}
{Aynı yol için hız \(\times\) zaman \(=\) yol sabit:
\[60\cdot 5=300\text{ km},\qquad 75\cdot t=300\ \Rightarrow\ t=4\text{ saat}\]}

\soru{Musluk -- Örnek 19}
{Eşit kapasiteli \(6\) musluk boş bir havuzu \(20\) saatte dolduruyor. Aynı türden \(8\) musluk havuzu kaç saatte doldurur?}
{Musluk sayısı \(\times\) süre sabit (havuz aynı):
\[6\cdot 20=8\cdot t\ \Rightarrow\ 120=8t\ \Rightarrow\ t=15\text{ saat}\]}

\soru{Ters Orantılı Paylaştırma -- Örnek 20}
{\(300\) TL, \(2\), \(3\) ve \(6\) sayıları ile \textbf{ters orantılı} olarak paylaştırılıyor. Her pay kaç TL'dir?}
{Paylar \(\dfrac{1}{2},\ \dfrac{1}{3},\ \dfrac{1}{6}\) ile doğru orantılı. Paydaların EKOK'u \(6\) ile çarpalım:
\[\tfrac{1}{2}:\tfrac{1}{3}:\tfrac{1}{6}=3:2:1\ \Rightarrow\ 6k=300\ \Rightarrow\ k=50\]
Paylar: \(150\), \(100\), \(50\) TL (en büyük sayıya en az pay).}

\subsection{Bileşik orantı}
\begin{frame}{Bileşik Orantı}
\begin{block}{Kural}
\(a\); \(b\) ile doğru, \(c\) ile ters orantılı ise
\[a=k\cdot\frac{b}{c}\qquad\text{yani}\qquad \frac{a\cdot c}{b}=k\ \text{(sabit)}\]
\end{block}
\begin{block}{İş problemleri (aynı verimle çalışanlar)}
\[\frac{\text{İş miktarı}}{\text{İşçi sayısı}\cdot\text{Günlük saat}\cdot\text{Gün}}=\text{sabit}\]
\[\frac{\text{İş}_1}{n_1\cdot s_1\cdot g_1}=\frac{\text{İş}_2}{n_2\cdot s_2\cdot g_2}\]
\end{block}
\end{frame}

\soru{Bileşik Orantı -- Örnek 21}
{\(x\) sayısı \(y\) ile doğru, \(z\) ile ters orantılıdır. \(x=4\), \(y=2\) iken \(z=6\)'dır. \(x=8\), \(y=6\) iken \(z\) kaçtır?}
{\(x=k\cdot\dfrac{y}{z}\ \Rightarrow\ k=\dfrac{x\cdot z}{y}=\dfrac{4\cdot 6}{2}=12\).
\[8=12\cdot\frac{6}{z}\ \Rightarrow\ 8z=72\ \Rightarrow\ z=9\]}

\soru{Bileşik Orantı -- Örnek 22}
{\(a\) sayısı \(b\) ile ters, \(c\) ile doğru orantılıdır. \(b=2\), \(c=3\) iken \(a=6\)'dır. \(b=4\), \(c=10\) iken \(a\) kaçtır?}
{\(a=k\cdot\dfrac{c}{b}\ \Rightarrow\ k=\dfrac{a\cdot b}{c}=\dfrac{6\cdot 2}{3}=4\).
\[a=4\cdot\frac{10}{4}=10\]}

\soru{İş Problemi -- Örnek 23}
{\(6\) işçi günde \(8\) saat çalışarak \(24\) m\(^2\) halıyı \(5\) günde dokuyor. \(10\) işçi günde \(12\) saat çalışarak \(36\) m\(^2\) halıyı kaç günde dokur?}
{\[\frac{24}{6\cdot 8\cdot 5}=\frac{36}{10\cdot 12\cdot x}\ \Rightarrow\ \frac{24}{240}=\frac{36}{120x}\]
\[\frac{1}{10}=\frac{36}{120x}\ \Rightarrow\ 120x=360\ \Rightarrow\ x=3\text{ gün}\]}

\soru{İş Problemi -- Örnek 24}
{\(4\) usta günde \(6\) saat çalışarak \(10\) günde \(40\) çift ayakkabı yapıyor. \(6\) usta günde \(8\) saat çalışarak \(5\) günde kaç çift ayakkabı yapar?}
{\[\frac{40}{4\cdot 6\cdot 10}=\frac{x}{6\cdot 8\cdot 5}\ \Rightarrow\ \frac{40}{240}=\frac{x}{240}\]
\[x=40\text{ çift}\]}

\begin{frame}{Ders 2 -- Sınıfta Siz Çözün}
\begin{columns}[t]
\column{0.62\textwidth}
\begin{enumerate}
\item \(3\) kg elma \(45\) TL ise \(7\) kg elma kaç TL?
\item \(8\) işçi bir işi \(12\) günde bitiriyor. Aynı işi \(6\) günde bitirmek için kaç işçi gerekir? \(\star\)
\item Eşit \(5\) musluk bir havuzu \(12\) saatte dolduruyor. \(3\) musluk kaç saatte doldurur?
\item \(x\); \(y\) ile doğru, \(z\) ile ters orantılı. \(x=6,\ y=4,\ z=2\). \ \(y=10,\ z=5\) iken \(x=?\) \(\star\)
\item \(5\) işçi günde \(8\) saatte \(10\) günde \(200\) m yol yapıyor. \(4\) işçi günde \(10\) saatte \(12\) günde kaç m yapar? \(\star\)
\end{enumerate}
\column{0.34\textwidth}
\pause
\textbf{Cevaplar:}
\begin{enumerate}
\item \(105\) TL
\item \(16\) işçi
\item \(20\) saat
\item \(6\)
\item \(240\) m
\end{enumerate}
\end{columns}
\vspace{4pt}{\footnotesize \(\star\) işaretli soruların adım adım çözümü sonraki slaytlardadır.}
\end{frame}

% ---- Ders 2 alıştırma çözümleri ----
\soru{Ders 2 -- Alıştırma Çözümü (Soru 2)}
{\(8\) işçi bir işi \(12\) günde bitiriyor. Aynı işi \(6\) günde bitirmek için aynı nitelikte kaç işçi gerekir?}
{Aynı iş için işçi sayısı ile gün ters orantılıdır:
\[8\cdot 12=n\cdot 6\ \Rightarrow\ 96=6n\ \Rightarrow\ n=16\text{ işçi}\]
Gün yarıya indi \(\Rightarrow\) işçi \(2\) katına çıktı \checkmark}

\soru{Ders 2 -- Alıştırma Çözümü (Soru 4)}
{\(x\); \(y\) ile doğru, \(z\) ile ters orantılıdır. \(x=6\), \(y=4\), \(z=2\) iken; \(y=10\), \(z=5\) olduğunda \(x\) kaçtır?}
{\(x=k\cdot\dfrac{y}{z}\ \Rightarrow\ k=\dfrac{x\cdot z}{y}=\dfrac{6\cdot 2}{4}=3\).
\[x=3\cdot\frac{10}{5}=6\]}

\soru{Ders 2 -- Alıştırma Çözümü (Soru 5)}
{\(5\) işçi günde \(8\) saat çalışarak \(10\) günde \(200\) m yol yapıyor. \(4\) işçi günde \(10\) saat çalışarak \(12\) günde kaç m yol yapar?}
{\[\frac{200}{5\cdot 8\cdot 10}=\frac{x}{4\cdot 10\cdot 12}\ \Rightarrow\ \frac{200}{400}=\frac{x}{480}\]
\[\frac{1}{2}=\frac{x}{480}\ \Rightarrow\ x=240\text{ m}\]}


% =====================================================
\section{Ders 3: Ortalamalar}
\begin{frame}
\centering\Large\color{anamavi}\textbf{Ders 3}\\[6pt]\normalsize Aritmetik, Ağırlıklı, Geometrik ve Harmonik Ortalama
\end{frame}

\subsection{Aritmetik ortalama}
\begin{frame}{Aritmetik Ortalama}
\begin{block}{Tanım}
\[\text{AO}=\frac{x_1+x_2+\cdots+x_n}{n}\qquad\Longrightarrow\qquad \text{Toplam}=n\cdot\text{AO}\]
Sayı eklenir/çıkarılırsa: \ yeni ortalama \(=\dfrac{\text{eski toplam}\pm\text{sayı}}{\text{yeni sayı adedi}}\)
\end{block}
\begin{block}{Ağırlıklı ortalama}
\(n_1\) kişilik grubun ortalaması \(o_1\), \(n_2\) kişilik grubunki \(o_2\) ise: \ \(\text{Ortalama}=\dfrac{n_1\,o_1+n_2\,o_2}{n_1+n_2}\)
\end{block}
\begin{alertblock}{Sık yapılan hatalar}
Farklı büyüklükteki grupların ortalaması \(\frac{o_1+o_2}{2}\) \textbf{değildir}; toplamlar üzerinden hesaplayın.
\end{alertblock}
\end{frame}

\soru{Aritmetik Ortalama -- Örnek 25}
{\(7\), \(12\), \(15\), \(26\) sayılarının aritmetik ortalaması kaçtır?}
{\[\text{AO}=\frac{7+12+15+26}{4}=\frac{60}{4}=15\]}

\soru{Aritmetik Ortalama -- Örnek 26}
{Bir öğrencinin ilk iki sınav notu \(60\) ve \(70\)'tir. Üç sınavın ortalamasının \(75\) olması için üçüncü sınavdan kaç alması gerekir?}
{Hedef toplam: \(3\cdot 75=225\). \quad Mevcut toplam: \(60+70=130\).
\[\text{Gereken not}=225-130=95\]}

\soru{Sayı Çıkarma -- Örnek 27}
{Aritmetik ortalaması \(15\) olan \(10\) sayıdan, toplamları \(30\) olan \(4\) sayı çıkarılıyor. Kalan sayıların aritmetik ortalaması kaçtır?}
{Toplam: \(10\cdot 15=150\). \quad Kalan toplam: \(150-30=120\).
\[\text{Kalan sayı adedi}=6,\qquad \text{AO}=\frac{120}{6}=20\]}

\soru{Sayı Ekleme -- Örnek 28}
{Yaş ortalaması \(24\) olan \(5\) kişilik bir gruba \(30\) yaşında bir kişi katılıyor. Grubun yeni yaş ortalaması kaçtır?}
{Eski toplam: \(5\cdot 24=120\). \quad Yeni toplam: \(120+30=150\).
\[\text{Yeni ortalama}=\frac{150}{6}=25\]}

\soru{Ağırlıklı Ortalama -- Örnek 29}
{\(20\) kişilik bir şubenin sınav ortalaması \(60\), \(30\) kişilik diğer şubenin ortalaması \(70\)'tir. İki şubenin birlikte ortalaması kaçtır?}
{\[\frac{20\cdot 60+30\cdot 70}{20+30}=\frac{1200+2100}{50}=\frac{3300}{50}=66\]
Dikkat: \(\frac{60+70}{2}=65\) yanlış; kalabalık şubeye daha yakın.}

\soru{Ağırlıklı Ortalama -- Örnek 30}
{Bir derste vize notunun ağırlığı \(2\), final notunun ağırlığı \(3\)'tür. Vizeden \(50\), finalden \(80\) alan öğrencinin ağırlıklı ortalaması kaçtır?}
{\[\frac{2\cdot 50+3\cdot 80}{2+3}=\frac{100+240}{5}=\frac{340}{5}=68\]}

\soru{Karışık Sınıf -- Örnek 31}
{Bir sınıfta kızların not ortalaması \(70\), erkeklerin \(60\), sınıfın ortalaması \(64\)'tür. Kız sayısının erkek sayısına oranı kaçtır?}
{Kız sayısı \(k\), erkek sayısı \(e\):
\[70k+60e=64(k+e)\ \Rightarrow\ 6k=4e\ \Rightarrow\ \frac{k}{e}=\frac{4}{6}=\frac{2}{3}\]
Kontrol: \(k=2,\ e=3\): \(\frac{140+180}{5}=64\) \checkmark}

\subsection{Geometrik ve harmonik ortalama}
\begin{frame}{Geometrik ve Harmonik Ortalama}
\begin{block}{İki pozitif sayı \(a\) ve \(b\) için}
\[\text{AO}=\frac{a+b}{2},\qquad \text{GO}=\sqrt{a\cdot b},\qquad \text{HO}=\frac{2ab}{a+b}\]
\[\text{AO}\ \ge\ \text{GO}\ \ge\ \text{HO}\qquad(\text{eşitlik yalnızca } a=b \text{ iken})\]
\end{block}
\begin{block}{Ortalama hız}
\[\text{Ortalama hız}=\frac{\text{Toplam yol}}{\text{Toplam zaman}}\]
Aynı yol \(v_1\) hızıyla gidilip \(v_2\) hızıyla dönülürse: \ \(\dfrac{2v_1v_2}{v_1+v_2}\) \ (harmonik ortalama).
\end{block}
\begin{alertblock}{Sık yapılan hatalar}
Gidiş--dönüşte ortalama hız \(\frac{v_1+v_2}{2}\) \textbf{değildir}.
\end{alertblock}
\end{frame}

\soru{Üç Ortalama -- Örnek 32}
{\(4\) ve \(9\) sayılarının aritmetik, geometrik ve harmonik ortalamalarını bulup sıralayınız.}
{\[\text{AO}=\frac{4+9}{2}=6{,}5,\qquad \text{GO}=\sqrt{36}=6,\qquad \text{HO}=\frac{2\cdot 36}{13}=\frac{72}{13}\approx 5{,}54\]
\[6{,}5>6>5{,}54\quad(\text{AO}>\text{GO}>\text{HO})\]}

\soru{Geometrik Ortalama -- Örnek 33}
{\(a\) ile \(b\)'nin geometrik ortalaması \(4\), aritmetik ortalaması \(5\)'tir. \(a^2+b^2\) kaçtır?}
{\(\sqrt{ab}=4\Rightarrow ab=16\); \quad \(\dfrac{a+b}{2}=5\Rightarrow a+b=10\).
\[a^2+b^2=(a+b)^2-2ab=100-32=68\]
(Sayılar: \(2\) ve \(8\); \ \(4+64=68\) \checkmark)}

\soru{Ortalama Hız -- Örnek 34}
{Bir araç A'dan B'ye saatte \(60\) km hızla gidip aynı yoldan saatte \(40\) km hızla dönüyor. Bütün yolculuktaki ortalama hızı kaç km/sa'tir?}
{\[\frac{2\cdot 60\cdot 40}{60+40}=\frac{4800}{100}=48\text{ km/sa}\]
Kontrol: yol \(120\) km olsun: \(2\) sa \(+\ 3\) sa \(=5\) sa; \ \(\dfrac{240}{5}=48\). \ (\(50\) değil!)}

\subsection{Ardışık sayılar}
\begin{frame}{Ardışık Sayıların Ortalaması}
\begin{block}{Kural}
Ardışık (eşit aralıklı) sayılarda
\[\text{Ortalama}=\frac{\text{ilk terim}+\text{son terim}}{2}\]
Terim sayısı tek ise bu, \textbf{ortanca} (tam ortadaki) terimdir.\\[2pt]
Toplam \(=\) terim sayısı \(\times\) ortalama.
\end{block}
\begin{alertblock}{Sık yapılan hatalar}
Terim sayısı: \(\dfrac{\text{son}-\text{ilk}}{\text{artış}}+1\) \quad (\(+1\)'i unutmayın).
\end{alertblock}
\end{frame}

\soru{Ardışık Sayılar -- Örnek 35}
{Ardışık \(5\) tek sayının toplamı \(85\)'tir. En büyüğü kaçtır? \ \(13\)'ten \(47\)'ye kadar (dahil) tek sayıların ortalaması ve toplamı kaçtır?}
{Ortanca terim \(=\dfrac{85}{5}=17\): \ \(13,15,\mathbf{17},19,21\) \(\Rightarrow\) en büyük \(21\).\\[3pt]
Ortalama \(=\dfrac{13+47}{2}=30\); \ terim sayısı \(=\dfrac{47-13}{2}+1=18\); \ toplam \(=18\cdot 30=540\).}

\begin{frame}{Ders 3 -- Sınıfta Siz Çözün}
\begin{columns}[t]
\column{0.62\textwidth}
\begin{enumerate}
\item \(8\) sayının ortalaması \(12\). Bu sayılara \(21\) eklenirse yeni ortalama? \(\star\)
\item \(3\) ve \(27\)'nin geometrik ortalaması?
\item Gidiş \(30\) km/sa, dönüş \(60\) km/sa (aynı yol). Ortalama hız? \(\star\)
\item \(15\) kişilik grubun ortalaması \(40\), \(10\) kişilik grubunki \(50\). Birlikte? \(\star\)
\item Ardışık \(7\) çift sayının toplamı \(98\). En küçüğü?
\end{enumerate}
\column{0.34\textwidth}
\pause
\textbf{Cevaplar:}
\begin{enumerate}
\item \(13\)
\item \(9\)
\item \(40\) km/sa
\item \(44\)
\item \(8\)
\end{enumerate}
\end{columns}
\vspace{4pt}{\footnotesize \(\star\) işaretli soruların adım adım çözümü sonraki slaytlardadır.}
\end{frame}

% ---- Ders 3 alıştırma çözümleri ----
\soru{Ders 3 -- Alıştırma Çözümü (Soru 1)}
{\(8\) sayının aritmetik ortalaması \(12\)'dir. Bu sayılara \(21\) eklenirse yeni ortalama kaç olur?}
{Eski toplam: \(8\cdot 12=96\). \quad Yeni toplam: \(96+21=117\), \ sayı adedi \(9\).
\[\text{Yeni ortalama}=\frac{117}{9}=13\]}

\soru{Ders 3 -- Alıştırma Çözümü (Soru 3)}
{Bir araç bir yolu saatte \(30\) km hızla gidip aynı yolu saatte \(60\) km hızla dönüyor. Ortalama hızı kaç km/sa'tir?}
{\[\frac{2\cdot 30\cdot 60}{30+60}=\frac{3600}{90}=40\text{ km/sa}\]
Kontrol: yol \(60\) km: \(2\) sa \(+\ 1\) sa \(=3\) sa; \ \(\dfrac{120}{3}=40\) \checkmark \ (\(45\) değil).}

\soru{Ders 3 -- Alıştırma Çözümü (Soru 4)}
{\(15\) kişilik bir grubun ortalaması \(40\), \(10\) kişilik bir grubun ortalaması \(50\)'dir. İki grubun birlikte ortalaması kaçtır?}
{Toplamlar: \(15\cdot 40=600\), \ \(10\cdot 50=500\).
\[\frac{600+500}{15+10}=\frac{1100}{25}=44\]
(\(\frac{40+50}{2}=45\) yanlış; kalabalık grup ortalamayı kendine çeker.)}


\begin{frame}{Özet}
\begin{columns}[t]
\column{0.32\textwidth}
\begin{block}{Ders 1}
\begin{itemize}
\item Oran birimsiz; birimler eşitlenir
\item \(\frac{a}{b}=\frac{c}{d}\Leftrightarrow ad=bc\)
\item \(\frac{a}{2}=\frac{b}{3}=\frac{c}{5}=k\): \(a=2k,\dots\)
\item Ölçek: harita / gerçek
\end{itemize}
\end{block}
\column{0.32\textwidth}
\begin{block}{Ders 2}
\begin{itemize}
\item Doğru: \(\frac{y}{x}=k\)
\item Ters: \(x\cdot y=k\)
\item Bileşik: \(a=k\frac{b}{c}\)
\item \(\frac{\text{İş}}{n\cdot s\cdot g}\) sabit
\end{itemize}
\end{block}
\column{0.32\textwidth}
\begin{block}{Ders 3}
\begin{itemize}
\item Toplam \(=n\cdot\)AO
\item Ağırlıklı: \(\frac{n_1o_1+n_2o_2}{n_1+n_2}\)
\item AO \(\ge\) GO \(\ge\) HO
\item Gidiş--dönüş: \(\frac{2v_1v_2}{v_1+v_2}\)
\end{itemize}
\end{block}
\end{columns}
\vspace{6pt}
\begin{alertblock}{Gelecek hafta}
Sayı, kesir ve yaş problemleri
\end{alertblock}
\end{frame}

\begin{frame}
\centering\Large İyi çalışmalar
\end{frame}

\end{document}
